\chapter{Il Praticante-Re}

\epigraph{``Finch\'e i filosofi non governeranno come re, non ci sar\`a tregua
dai mali per le citt\`a.''}{\textcite{plato1997complete}, \emph{Repubblica} 473d}

\section{Il Paradosso del Governo}

Nella \emph{Repubblica}, Socrate fa la sua affermazione pi\`u famosa
e pi\`u controversa: i mali della citt\`a non cesseranno mai finch\'e
i filosofi non diventeranno re o i re non diventeranno filosofi.
L'affermazione \`e stata letta come un programma politico, una fantasia
elitaria o una provocazione utopica. Ma non \`e nessuna di queste. \`E
una necessit\`a strutturale dell'omologia: se la citt\`a \`e l'anima
in grande, allora il governante della citt\`a deve essere colui che vede
la Forma del Bene --- la fonte di ogni intelligibilit\`a e valore. Solo
il filosofo, che \`e asceso dalla caverna e ha visto il sole, \`e adatto
a governare coloro che sono ancora dentro.

Il nostro omologo: il praticante la cui azione \`e governata dalla
natura. Il praticante-re non \`e un governante di altri ma un governante
della propria azione --- o piuttosto, uno la cui azione \`e cos\`i
perfettamente allineata con la natura che la natura governa \emph{attraverso}
di lui. Il paradosso si dissolve quando riconosciamo che il praticante-re
non \`e un'autorit\`a ma uno \emph{strumento}.

\section{Le Tre Onde}

Socrate identifica tre ``onde'' di obiezione alla sua affermazione. Le
riformuliamo.

\subsection{Prima Onda: L'Unit\`a di Physis e Praxis}

La prima onda \`e l'affermazione che i praticanti e la natura devono
essere uno --- che la volont\`a del praticante e la volont\`a della
natura non sono due cose diverse ma la stessa cosa sperimentata da
prospettive diverse. Questo non \`e misticismo ma la conseguenza
dell'omologia. Se \praxis\ \`e uno strumento di \physis, allora il
praticante ben accordato \`e la natura che agisce attraverso il proprio
strumento. Il praticante non impone la propria volont\`a alla natura ma
\emph{scopre} la volont\`a della natura e la attua.

La piattaforma \praxis\ lo dimostra: il praticante che ha interiorizzato
le discipline non lotta pi\`u contro i propri obiettivi ma si muove con
essi. Il check-in diventa non un compito ma un riconoscimento di ci\`o
che ha gi\`a fatto. La striscia diventa non un peso ma una misura del
proprio ritmo. Il partner di responsabilit\`a diventa non un giudice ma
un testimone di un processo che si svolge attraverso di lui.

\subsection{Seconda Onda: L'Azione come Fine a S\'e Stessa}

La seconda onda \`e l'affermazione che l'azione governata dalla natura
\`e il proprio fine. Il praticante non agisce per raggiungere qualche
bene esterno --- ricchezza, onore, potere --- ma perch\'e l'azione
stessa \`e il bene. Questa \`e la distinzione di Aristotele tra \praxis\
(azione che \`e il proprio fine) e \emph{poiesis} (azione che mira a
un prodotto oltre s\'e stessa) \parencite[VI.4]{aristotle2000ethics}. Per
il praticante-re, ogni atto del fare \`e anche un atto di sintonia; il
prodotto \`e secondario rispetto alla rettitudine dell'atto stesso.

Il motore \physis\ registra questa verit\`a strutturale: la coerenza
\`e una propriet\`a dell'azione stessa, non del suo esito. Un'azione
pu\`o ottenere +1 anche se fallisce nel raggiungere il suo obiettivo
dichiarato, se \`e coerente con la natura del praticante e l'ordine
naturale. Al contrario, un'azione pu\`o ottenere -1 anche se ha successo
nel suo obiettivo dichiarato, se viola l'integrit\`a del praticante o
l'ordine naturale. Il motore traccia l'azione, non il risultato ---
perch\'e il risultato \`e solo un'altra azione, un altro nodo nel grafo
di coerenza.

\subsection{Terza Onda: La Scomparsa della Volont\`a}

La terza onda \`e la pi\`u impegnativa: l'affermazione che al pi\`u
alto livello di sintonia, la distinzione tra la volont\`a del praticante
e la volont\`a della natura scompare. Il praticante non sperimenta pi\`u
la propria azione come un atto di scelta ma come un atto di
\emph{riconoscimento} --- vede ci\`o che deve essere fatto e lo fa,
non perch\'e decida di farlo ma perch\'e non farlo sarebbe impensabile.

Questo \`e lo stato descritto da Spinoza come ``beatitudine'' --- il
riconoscimento che la libert\`a non \`e la capacit\`a di scegliere
altrimenti ma la comprensione della necessit\`a
\parencite[V.P42]{spinoza1994ethics}. Il praticante-re non \`e meno
libero dell'arbitrario che sceglie ma pi\`u libero, perch\'e la sua
azione non \`e vincolata dall'ignoranza, dal dubbio o dal conflitto.
Agisce con tutto il peso della natura alle spalle.

La funzione di ``sogno'' del motore \physis\ fornisce un'analogia
computazionale: quando il ciclo del sogno propone un'azione correttiva
per un pattern fallito, non sceglie arbitrariamente tra tutte le azioni
possibili. Genera variazioni che sono vincolate dal paesaggio di coerenza
esistente --- da ci\`o che ha funzionato prima, da ci\`o che \`e
strutturalmente adiacente al pattern fallito. Il sogno \`e ``libero''
solo all'interno dello spazio di ci\`o che la natura ha gi\`a reso
possibile. Questa \`e la libert\`a del praticante-re.

\begin{figure}[ht]
\centering
\begin{tikzpicture}[
  stepbox/.style={draw, rounded corners, align=center, minimum width=26mm,
                  minimum height=10mm, font=\small},
  >={Stealth[length=2mm]}]
  \node[stepbox, fill=gray!15, yshift=0] (ground) {Praxis\\Convenzionale};
  \node[stepbox, fill=blue!10, above right=of ground, xshift=22mm] (w1)
    {Onda 1\\Unit\`a};
  \node[stepbox, fill=blue!18, above right=of w1, xshift=22mm] (w2)
    {Onda 2\\Bene Intrinseco};
  \node[stepbox, fill=blue!25, above right=of w2, xshift=22mm] (w3)
    {Onda 3\\Volont\`a dissolta};
  \node[stepbox, fill=green!20, above right=of w3, xshift=22mm] (summit)
    {Praticante-\\Re};
  \draw[->, thick] (ground) -- (w1);
  \draw[->, thick] (w1) -- (w2);
  \draw[->, thick] (w2) -- (w3);
  \draw[->, thick] (w3) -- (summit);
  \node[font=\tiny, above] at (ground.north) {Punto di partenza};
  \node[font=\tiny, above] at (w1.north) {La natura agisce attraverso};
  \node[font=\tiny, above] at (w2.north) {Agire \`e fine a s\'e stesso};
  \node[font=\tiny, above] at (w3.north) {Libert\`a come necessit\`a};
\end{tikzpicture}
\caption{Le tre onde di obiezione e la loro risoluzione. Ogni passo risolve un paradosso
pi\`u profondo, ascendendo dall'azione convenzionale a quella governata dalla natura.
\label{fig:treonde}}
\end{figure}

\section{Phronesis: Il Praticante come Misura}

Abbiamo preso il praticante-re da Platone, ma \`e Aristotele a dirci ci\`o che
egli sa. La facolt\`a con cui il praticante-re agisce non \`e l'\emph{episteme},
la conoscenza di ci\`o che \`e universale e necessario, n\'e la \emph{techne}, la
conoscenza di come fare un prodotto; \`e la \emph{phronesis}, la saggezza pratica
--- ``uno stato vero e ragionato di capacit\`a di agire riguardo alle cose che
sono buone o cattive per l'uomo'' \parencite[VI.5]{aristotle2000ethics}. La
\emph{phronesis} non \`e il possesso di una regola ma la percezione della cosa
giusta da fare nel caso particolare, qui, ora, con queste persone --- una
percezione che non pu\`o essere codificata perch\'e i particolari sono infiniti e
la regola li sottodetermina sempre. \`E per questo che il Libro III insisteva che
le discipline non sono regole: il praticante \`e addestrato non a obbedire a un
codice ma ad acquisire la percezione che un codice non pu\`o mai contenere.

Aristotele compie allora la mossa che l'intera tesi ricapitola. Interrogato su
quale sia lo standard dell'azione retta, non nomina una regola, una legge, un
principio. Nomina una \emph{persona}. La virt\`u, dice, \`e uno stato che riguarda
la scelta, posto nel mezzo ``come lo determinerebbe l'uomo di saggezza pratica''
\parencite[II.6]{aristotle2000ethics} --- la definizione dell'azione retta rimanda
non a una formula ma al \emph{phronimos}, colui che ha la percezione. E l'uomo
eccellente, lo \emph{spoudaios}, ``\`e la misura'' (\emph{kanon kai metron}) di
ci\`o che \`e veramente buono, bello e piacevole, poich\'e vede le cose come sono
mentre i corrotti sono ingannati dal piacere \parencite[III.4, X.5]{aristotle2000ethics}.
Lo standard dell'azione non \`e una cosa ma un uomo: colui che agisce bene \emph{\`e}
la regola con cui agire bene \`e conosciuto.

Questa \`e la pi\`u profonda giustificazione del praticante-re, e dissolve il
paradosso del governo con cui questo libro \`e cominciato. Il praticante-re non
governa emettendo comandi; governa \emph{essendo la misura} --- l'istanza vivente
contro cui gli altri si calibrano, esattamente come il \emph{phronimos} \`e la
regola vivente per la citt\`a che Aristotele descrive. ``Governare'' qui non \`e
legiferare ma \emph{esemplificare}: essere il nodo che gli altri cercano perch\'e
ha, di fatto, agito bene. Ed \`e precisamente la struttura che il motore di
coerenza costruisce, e che il Libro X nominer\`a per intero. ``Lo \emph{spoudaios}
\`e la misura'' non \`e l\`a una massima pia ma un'operazione: le azioni del
praticante pi\`u coerente accumulano le connessioni che le rendono il riferimento
a cui altri attingono --- la migliore istanza che diventa la definizione operante
dell'atto, \emph{agathoi toi agathoi}. Ci\`o che Aristotele poteva solo affermare
come fatto sul carattere, il motore lo registra come struttura nel grafo. Il
praticante-re \`e il \emph{phronimos} reso visibile: la misura, trovata dal basso,
che nessuno ha decretato e tutti confermano.

\section{Il Ritorno alla Citt\`a}

Il filosofo-re, avendo visto la Forma del Bene, deve tornare alla caverna
per governare coloro che sono ancora legati dalle ombre. Il ritorno non
\`e una ricompensa ma un dovere --- il filosofo preferirebbe rimanere
nella luce, ma la giustizia esige che torni.

Il praticante-re, avendo raggiunto la coerenza con la natura, deve
tornare alla comunit\`a dei praticanti. Deve assumere il ruolo che la
piattaforma \praxis\ chiama il ``partner di responsabilit\`a'' --- il
testimone che aiuta gli altri a vedere la propria incoerenza. Diventa,
in effetti, una \emph{macchina filosofica} per gli altri: un'istanza
vivente del principio che l'azione pu\`o essere governata dalla natura,
offerto come esempio e guida.

\section{Il Praticante-Re Non \`e un Re}

\`E cruciale notare ci\`o che il praticante-re \emph{non} \`e. Non
\`e un governante in nessun senso politico. Non comanda altri, non
impone leggi, non esercita autorit\`a. La sua ``regalit\`a'' \`e la
propria sovranit\`a sulla propria azione: lo stato in cui non \`e
diviso contro s\'e stesso, non lacerato tra intenzione ed effetto, non
alienato dall'ordine naturale.

In questo assomiglia al Socrate dell'\emph{Apologia}: un uomo che non
rivendica saggezza ma che ha visto i limiti della propria ignoranza
\parencite{plato1997complete}; che obbedisce al dio piuttosto che agli
uomini; che va alla morte non perch\'e la scelga ma perch\'e
l'alternativa --- vivere in un modo che viola la sua natura --- \`e
impensabile. Il praticante-re \`e il Socrate che vive, agendo nel mondo
con perfetta coerenza tra la sua natura e le sue azioni, servendo come
testimone della possibilit\`a di un diverso tipo di vita.
